
% UART Spec  (as_uart_top / ip_uart_v2)
\label{sec:uartfunc01}

%=============================================================
% Requirements
%=============================================================
\subsection{Requirements}
\label{subsec:uartreq01}

This section defines the functional requirements for the \textbf{UART0}
peripheral (module \texttt{as\_uart\_top}) in the \aschip{} system.
Requirements follow IEEE\;29148 terminology:
\textit{shall} = mandatory, \textit{should} = recommended.

Each requirement identifier has the form
\textbf{UART-\textit{CAT}-\textit{NNN}}, where \textit{CAT} is the
category code (see Table~\ref{tab:asUartReqCat}) and \textit{NNN} is a
three-digit sequence number.

\begin{table}[H]
\caption{UART Requirements -- Category Codes}
\label{tab:asUartReqCat}
\centering
\begin{tabularx}{\textwidth}{|l |X|}
  \hline
  Category Code & Scope \\
  \hline
  \hline
  FR & Frame Format \\
  \hline
  BR & Baud Rate and Oversampling \\
  \hline
  TX & Transmitter Path \\
  \hline
  RX & Receiver Path \\
  \hline
  FI & FIFOs \\
  \hline
  ER & Error Detection \\
  \hline
  IR & Interrupts \\
  \hline
  SW & Software Interface (Register Map) \\
  \hline
  HW & Hardware Interface (Pins) \\
  \hline
\end{tabularx}
\end{table}

%--------------------------------------------------------------
\subsubsection{Frame Format Requirements (FR)}
%--------------------------------------------------------------

\begin{table}[H]
\caption{UART Requirements -- Frame Format}
\label{tab:asUartReqFR}
\centering
\begin{tabularx}{\textwidth}{|l |X |X|}
  \hline
  ID & Requirement & Notes / Rationale \\
  \hline
  \hline
  UART-FR-001 &
    The UART \textit{shall} support data word lengths of 5, 6, 7, and
    8\;bits per frame, configurable via \textbf{UART0\_LCR}[1:0]:
    00\,=\,5\;bit, 01\,=\,6\;bit, 10\,=\,7\;bit, 11\,=\,8\;bit (reset). &
    Covers the full range of the standard asynchronous serial frame (RS-232 / UART). \\
  \hline
  UART-FR-002 &
    The UART \textit{shall} support three parity modes configurable via
    \textbf{UART0\_LCR}[3:2]:
    00\,=\,none (reset), 01\,=\,odd, 10\,=\,even.
    Encoding 11 is reserved (no parity). &
    Standard 8N1, 7E1, 7O1, etc.\ coverage.
    The parity bit is inserted between the last data bit and the first stop bit. \\
  \hline
  UART-FR-003 &
    The UART \textit{shall} support 1 stop bit (reset) and 2 stop bits,
    selectable via \textbf{UART0\_LCR}[4]:
    0\,=\,1 stop bit, 1\,=\,2 stop bits. &
    1.5 stop bits are omitted; they are only defined for 5-bit data frames
    and not required in this application. \\
  \hline
  UART-FR-004 &
    The UART \textit{shall} support forced break transmission:
    when \textbf{UART0\_LCR}[5]=1 the TX output \textit{shall} be driven
    continuously to logic\;0, independent of FIFO content and frame
    boundaries. &
    Required by POSIX \texttt{tcsendbreak()} and Linux \texttt{termios}.
    Break ends when LCR[5] is cleared. \\
  \hline
\end{tabularx}
\end{table}

%--------------------------------------------------------------
\subsubsection{Baud Rate and Oversampling Requirements (BR)}
%--------------------------------------------------------------

\begin{table}[H]
\caption{UART Requirements -- Baud Rate and Oversampling}
\label{tab:asUartReqBR}
\centering
\begin{tabularx}{\textwidth}{|l |X |X|}
  \hline
  ID & Requirement & Notes / Rationale \\
  \hline
  \hline
  UART-BR-001 &
    The UART \textit{shall} derive the baud-rate clock from the system clock
    $f_{clk}$ using a 16-bit integer divisor stored in
    \textbf{UART0\_CLKDIV}[15:0]:
    \[ f_{baud} = \frac{f_{clk}}{16 \times \text{CLKDIV}} \]
    CLKDIV\,=\,0 is illegal; the reset value \textit{shall} correspond to
    115\,200\;Baud at 125\;MHz (CLKDIV\,=\,68). &
    16$\times$ oversampling is the industry standard (16550 compatible).
    At 125\;MHz and CLKDIV\,=\,68 the actual baud rate is
    $\approx$115\,042\;Baud ($-0.14\,\%$ error). \\
  \hline
  UART-BR-002 &
    The valid range of CLKDIV \textit{shall} be 1\ldots65535.
    Writing CLKDIV\,=\,0 \textit{shall} have no effect (value rejected,
    previous value retained). &
    CLKDIV\,=\,0 would cause division-by-zero in the baud-rate counter. \\
  \hline
  UART-BR-003 &
    The oversampling ratio \textit{shall} be 16.
    Each received bit \textit{shall} be sampled at sub-clock\;8 of its
    16-sub-clock period (centre-of-bit sampling). &
    Centre sampling maximises noise margin relative to both edges. \\
  \hline
\end{tabularx}
\end{table}

%--------------------------------------------------------------
\subsubsection{Transmitter Requirements (TX)}
%--------------------------------------------------------------

\begin{table}[H]
\caption{UART Requirements -- Transmitter}
\label{tab:asUartReqTX}
\centering
\begin{tabularx}{\textwidth}{|l |X |X|}
  \hline
  ID & Requirement & Notes / Rationale \\
  \hline
  \hline
  UART-TX-001 &
    The transmitter \textit{shall} serialize data frames from the TX FIFO:
    start bit (0), data bits LSB first, parity bit (if enabled), stop
    bit(s) (1). &
    Standard UART bit ordering. \\
  \hline
  UART-TX-002 &
    The TX output (\texttt{tx\_o}) \textit{shall} be logic\;1 (mark state)
    when idle and between frames. &
    Required so that the receiver can detect the start-bit falling edge. \\
  \hline
  UART-TX-003 &
    The transmitter \textit{shall} fetch the next frame from the TX FIFO
    only when the FIFO is non-empty and the shift register is idle.
    If the TX FIFO is empty while the shift register finishes the current
    frame, \texttt{tx\_o} \textit{shall} return to idle (logic\;1). &
    Prevents underrun: partial frames are never transmitted. \\
  \hline
  UART-TX-004 &
    When \textbf{UART0\_LCR}[5]\,=\,1 (break), the transmitter
    \textit{shall} drive \texttt{tx\_o}\,=\,0 continuously until
    LCR[5] is cleared, overriding normal FIFO operation. &
    See UART-FR-004. \\
  \hline
\end{tabularx}
\end{table}

%--------------------------------------------------------------
\subsubsection{Receiver Requirements (RX)}
%--------------------------------------------------------------

\begin{table}[H]
\caption{UART Requirements -- Receiver}
\label{tab:asUartReqRX}
\centering
\begin{tabularx}{\textwidth}{|l |X |X|}
  \hline
  ID & Requirement & Notes / Rationale \\
  \hline
  \hline
  UART-RX-001 &
    The \texttt{rx\_i} input \textit{shall} be synchronised through a
    2-stage flip-flop synchroniser operating in the peripheral clock domain
    before any logic (edge detection, sampling) is applied to it. &
    Mandatory metastability protection for any asynchronous input on an
    FPGA; omitting it can cause the entire system to malfunction. \\
  \hline
  UART-RX-002 &
    Start-bit detection \textit{shall} be performed by detecting a
    high-to-low transition (falling edge) on the synchronised \texttt{rx\_i}
    signal. &
    Marks the beginning of a new frame. \\
  \hline
  UART-RX-003 &
    After start-bit detection, each data bit and (if enabled) the parity bit
    \textit{shall} be sampled at sub-clock\;8 of the corresponding bit
    period (centre sampling, see UART-BR-003). &
    The start-bit centre is used to align the oversampling counter so that
    subsequent bits are sampled at their centres. \\
  \hline
  UART-RX-004 &
    On successful frame completion (valid stop bit detected), the received
    data word \textit{shall} be written into the RX FIFO. &
    Frames with a missing stop bit generate a frame error (see UART-ER-001)
    and the data word is still pushed to the FIFO
    so that SW can inspect the erroneous byte. \\
  \hline
  UART-RX-005 &
    When \textbf{UART0\_CTRL}[0]\,=\,1 (loopback), the receiver input
    \textit{shall} be internally connected to the TX shift-register output,
    bypassing the external \texttt{rx\_i} pin. &
    Enables SW self-test without external hardware. \\
  \hline
\end{tabularx}
\end{table}

%--------------------------------------------------------------
\subsubsection{FIFO Requirements (FI)}
%--------------------------------------------------------------

\begin{table}[H]
\caption{UART Requirements -- FIFOs}
\label{tab:asUartReqFI}
\centering
\begin{tabularx}{\textwidth}{|l |X |X|}
  \hline
  ID & Requirement & Notes / Rationale \\
  \hline
  \hline
  UART-FI-001 &
    The peripheral \textit{shall} provide a TX FIFO and an RX FIFO.
    FIFO depth \textit{shall} be set by the SystemVerilog parameter
    \texttt{FIFO\_DEPTH} (default: 16 entries, 8\;bits per entry). &
    Decouples bus transaction timing from baud-rate timing; prevents
    CPU stalls on every byte. \\
  \hline
  UART-FI-002 &
    Each FIFO \textit{shall} provide the following status signals to the
    register block: full, empty, half-full (level\,>\,FIFO\_DEPTH/2),
    and current fill level as a $\lceil\log_2(\text{FIFO\_DEPTH}+1)\rceil$-bit
    integer. &
    Required for FIFOSTAT register (UART-SW-001). \\
  \hline
  UART-FI-003 &
    FIFO status \textit{shall} be readable via \textbf{UART0\_FIFOSTAT}. &
    Allows software to perform polled I/O without relying on interrupts. \\
  \hline
  UART-FI-004 &
    Writing 1 to \textbf{UART0\_CTRL}[1] (tx\_flush) \textit{shall}
    synchronously reset the TX FIFO to the empty state.
    The bit \textit{shall} auto-clear on the next clock cycle. &
    Used to abort a pending transmission without a full peripheral reset. \\
  \hline
  UART-FI-005 &
    Writing 1 to \textbf{UART0\_CTRL}[2] (rx\_flush) \textit{shall}
    synchronously reset the RX FIFO to the empty state.
    The bit \textit{shall} auto-clear on the next clock cycle. &
    Used to discard stale receive data after a configuration change. \\
  \hline
  UART-FI-006 &
    The RX interrupt threshold \textit{shall} be configurable via
    \textbf{UART0\_RXTHRES}[4:0] (valid range: 1\ldots FIFO\_DEPTH;
    reset: FIFO\_DEPTH/2\,=\,8).
    The rx\_ready interrupt source (RIS[0]) \textit{shall} be asserted when
    the RX FIFO level is $\geq$ RXTHRES. &
    Allows SW to tune the interrupt rate vs.\ latency trade-off without
    code changes in the driver. \\
  \hline
\end{tabularx}
\end{table}

%--------------------------------------------------------------
\subsubsection{Error Detection Requirements (ER)}
%--------------------------------------------------------------

\begin{table}[H]
\caption{UART Requirements -- Error Detection}
\label{tab:asUartReqER}
\centering
\begin{tabularx}{\textwidth}{|l |X |X|}
  \hline
  ID & Requirement & Notes / Rationale \\
  \hline
  \hline
  UART-ER-001 &
    The UART \textit{shall} detect a \textbf{frame error}:
    the expected stop bit is sampled as logic\;0.
    On detection, STATUS[4] (frame\_err) \textit{shall} be set. &
    Indicates baud-rate mismatch or line noise. \\
  \hline
  UART-ER-002 &
    The UART \textit{shall} detect a \textbf{parity error} when parity is
    enabled (LCR[3:2]\,$\neq$\,00): the received parity bit differs from
    the locally computed parity of the received data word.
    On detection, STATUS[5] (parity\_err) \textit{shall} be set. &
    Indicates single-bit data corruption. \\
  \hline
  UART-ER-003 &
    The UART \textit{shall} detect an \textbf{overrun error}:
    the RX FIFO is full when the receiver completes a new frame.
    The new byte \textit{shall} be discarded, and STATUS[6]
    (overrun\_err) \textit{shall} be set. &
    Signals that the driver is not reading fast enough. \\
  \hline
  UART-ER-004 &
    The UART \textit{shall} detect a \textbf{break condition}:
    \texttt{rx\_i} (synchronised) is held continuously at logic\;0 for
    $\geq$\;1 complete character time (start + data + parity + stop bits)
    after the initial falling edge.
    On detection, STATUS[7] (break\_det) \textit{shall} be set. &
    Required for Linux \texttt{tty} break handling and RS-232 modem
    signalling. \\
  \hline
  UART-ER-005 &
    The error flags STATUS[7:4] \textit{shall} be sticky (latched):
    they \textit{shall} remain set until explicitly cleared by writing\;1
    to the corresponding bit in \textbf{UART0\_ICR}. &
    Prevents loss of error information if the ISR runs after a delay. \\
  \hline
  UART-ER-006 &
    Each of the four error types (frame, parity, overrun, break)
    \textit{shall} generate an individual interrupt source in RIS
    (see UART-IR-004). &
    Allows driver to classify the error without reading STATUS separately. \\
  \hline
\end{tabularx}
\end{table}

%--------------------------------------------------------------
\subsubsection{Interrupt Requirements (IR)}
%--------------------------------------------------------------

\begin{table}[H]
\caption{UART Requirements -- Interrupts}
\label{tab:asUartReqIR}
\centering
\begin{tabularx}{\textwidth}{|l |X |X|}
  \hline
  ID & Requirement & Notes / Rationale \\
  \hline
  \hline
  UART-IR-001 &
    The UART \textit{shall} provide an \textbf{rx\_ready} interrupt source
    (RIS[0]) asserted when the RX FIFO fill level is $\geq$
    \textbf{UART0\_RXTHRES}. &
    Primary trigger for the receive interrupt service routine. \\
  \hline
  UART-IR-002 &
    The UART \textit{shall} provide a \textbf{tx\_ready} interrupt source
    (RIS[1]) asserted when the TX FIFO fill level falls below
    \textbf{UART0\_RXTHRES} (space available for writing). &
    Allows SW to refill the TX FIFO in an interrupt-driven transmit loop
    without polling. \\
  \hline
  UART-IR-003 &
    The UART \textit{shall} provide an \textbf{rx\_timeout} interrupt source
    (RIS[2]) asserted when the RX FIFO contains at least one byte and no
    new byte has been received for $\geq$\;4 character times. &
    Prevents data stranded in the RX FIFO below the threshold from being
    silently ignored when a frame does not exactly fill the FIFO to RXTHRES. \\
  \hline
  UART-IR-004 &
    The UART \textit{shall} provide individual interrupt sources:
    frame\_err (RIS[3]), parity\_err (RIS[4]),
    overrun\_err (RIS[5]), break\_det (RIS[6]). &
    See UART-ER-006. \\
  \hline
  UART-IR-005 &
    Each interrupt source \textit{shall} be individually maskable via
    the corresponding bit in \textbf{UART0\_IMSC}.
    The interrupt output signal (\texttt{uart\_irq\_o})
    \textit{shall} equal the bitwise OR of all (RIS\,\&\,IMSC) bits:
    \texttt{uart\_irq\_o}\,=\,\texttt{|}\,(RIS\,\&\,IMSC). &
    One-IRQ-line model consistent with GPIO0 and QSPI0. \\
  \hline
  UART-IR-006 &
    Any bit in RIS \textit{shall} be clearable by writing\;1 to the
    corresponding bit of \textbf{UART0\_ICR} (write-1-to-clear). &
    Consistent with QSPI0 interrupt model. \\
  \hline
\end{tabularx}
\end{table}

%--------------------------------------------------------------
\subsubsection{Software Interface Requirements (SW)}
%--------------------------------------------------------------

\begin{table}[H]
\caption{UART Requirements -- Software Interface}
\label{tab:asUartReqSW}
\centering
\begin{tabularx}{\textwidth}{|l |X |X|}
  \hline
  ID & Requirement & Notes / Rationale \\
  \hline
  \hline
  UART-SW-001 &
    All configuration, status, FIFO access, and interrupt control
    \textit{shall} be accessible through exactly 12 memory-mapped,
    64-bit registers on a Wishbone B4 Classic slave interface
    implemented by \texttt{as\_slave\_bpi}. &
    Consistent with GPIO0 and QSPI0. \\
  \hline
  UART-SW-002 &
    Registers \textit{shall} be 8-byte aligned and occupy the address
    range UART0\_BASE\,+\,0x00 to UART0\_BASE\,+\,0x5F (12\,$\times$\,8\;bytes). &
    Fits within the reserved address range
    0x0000\-0001\-0000\-0120 -- 0x0000\-0001\-0000\-017F. \\
  \hline
  UART-SW-003 &
    The base address of UART0 \textit{shall} be
    \texttt{0x0000\_0001\_0000\_0120}. &
    As defined in the system memory map (\texttt{as\_pack.sv}:
    \texttt{uart0\_start\_address\_c}). \\
  \hline
  UART-SW-004 &
    Write accesses to read-only register fields \textit{shall} have no
    effect; the field value \textit{shall} not change. &
    Standard HW register protection. \\
  \hline
  UART-SW-005 &
    Read accesses to write-only registers (ICR, and DATA when read
    returns RX FIFO data) \textit{shall} return the correct value:
    ICR reads as 0x0; DATA read pops the RX FIFO. &
    DATA register is bi-directional: write pushes TX FIFO, read pops
    RX FIFO (standard UART DATA register model). \\
  \hline
  UART-SW-006 &
    The Wishbone interface \textit{shall} respond with zero wait states
    (ACK asserted one clock after STB\,\&\,CYC). &
    Consistent with all other Wishbone slaves in the system. \\
  \hline
\end{tabularx}
\end{table}

%--------------------------------------------------------------
\subsubsection{Hardware Interface Requirements (HW)}
%--------------------------------------------------------------

\begin{table}[H]
\caption{UART Requirements -- Hardware Interface}
\label{tab:asUartReqHW}
\centering
\begin{tabularx}{\textwidth}{|l |X |X|}
  \hline
  ID & Requirement & Notes / Rationale \\
  \hline
  \hline
  UART-HW-001 &
    The peripheral \textit{shall} expose exactly two external serial I/O
    signals: \texttt{tx\_o} (output) and \texttt{rx\_i} (input). &
    Standard asynchronous UART interface; no modem control signals
    (RTS/CTS) required for this application. \\
  \hline
  UART-HW-002 &
    \texttt{rx\_i} \textit{shall} be treated as an asynchronous input
    and synchronised by a 2-stage flip-flop synchroniser in the
    peripheral clock domain before use (see UART-RX-001). &
    Mandatory for FPGA synthesis tool to infer proper placement
    of the input register chain and to suppress metastability. \\
  \hline
  UART-HW-003 &
    The peripheral clock (\texttt{clk\_i}) \textit{shall} be 125\;MHz
    (Zybo board system clock). &
    Determines the achievable baud rate resolution (UART-BR-001). \\
  \hline
\end{tabularx}
\end{table}

%=============================================================
% System Integration of UART0
%=============================================================
\subsection{System Integration of UART0}

%------------------------
\subsubsection{History}

\begin{table}[H]
\caption{UART0 Authors}
\label{tab:uartauth01}
\centering
\begin{tabularx}{\textwidth}{|X |X |X|}
  \hline
  Author & Department & From \\
  \hline
  \hline
  Andreas Siggelkow & Fac. E & June 2026 \\
  \hline
  & & \\
  \hline
\end{tabularx}
\end{table}

\begin{table}[H]
\caption{UART0 Section History}
\label{tab:uarthistory01}
\centering
\begin{tabularx}{\textwidth}{|X |X |X|}
  \hline
  Date & Name & Content Changes \\
  \hline
  \multicolumn{3}{|l|}{Version of this document: V0.1} \\
  \hline
  2026-06-05 & A. Siggelkow, Claude & First draft: Requirements section \\
  \hline
\end{tabularx}
\end{table}

%------------------------
\subsubsection{Functional Configuration of UART0}
\textcolor{red}{This section should describe the product-specific instantiation of this module's feature set.}

The standard features of the UART peripheral are described in its functional
description in Section~\ref{subsec:uartfunc01}.
Generic features for UART0 are listed in
Table~\ref{tab:asUartPer01a}.

\begin{table}[H]
\caption{UART0 Feature Set}
\label{tab:asUartPer01a}
\centering
\begin{tabularx}{\textwidth}{|l |l |X|}
  \hline
   & Generic UART Feature Set & Details \\
  \hline
  \hline
  yes & Separate Kernel Clock         & 125\;MHz \\
  \hline
  yes & Configurable Data Bits        & 5, 6, 7, 8 bits via LCR \\
  \hline
  yes & Configurable Parity           & None (reset), odd, even via LCR \\
  \hline
  yes & Configurable Stop Bits        & 1 (reset) or 2 via LCR \\
  \hline
  yes & Programmable Baud Rate        & $f_{baud} = 125\,\text{MHz} / (16 \times \text{CLKDIV})$; reset: 115\,200\;Baud \\
  \hline
  yes & TX FIFO                       & 16 entries $\times$ 8\;bit \\
  \hline
  yes & RX FIFO                       & 16 entries $\times$ 8\;bit \\
  \hline
  yes & RX FIFO Threshold             & Configurable 1\ldots16, reset: 8 \\
  \hline
  yes & Loopback Mode                 & CTRL[0] \\
  \hline
  yes & Break Generation              & LCR[5] \\
  \hline
  yes & RX Input Synchroniser         & 2-stage FF, mandatory metastability protection \\
  \hline
  yes & Frame Error Detection         & STATUS[4], RIS[3] \\
  \hline
  yes & Parity Error Detection        & STATUS[5], RIS[4] \\
  \hline
  yes & Overrun Error Detection       & STATUS[6], RIS[5] \\
  \hline
  yes & Break Detect                  & STATUS[7], RIS[6] \\
  \hline
  yes & RX Timeout Interrupt          & RIS[2], after 4 character times of FIFO inactivity \\
  \hline
  yes & Interrupt Generation          & 7 sources, maskable via IMSC \\
  \hline
  yes & FIFO Flush                    & tx\_flush / rx\_flush in CTRL \\
  \hline
  no  & 9-bit Data Words              & Not required \\
  \hline
  no  & 1.5 Stop Bits                 & Not required \\
  \hline
  no  & Hardware Flow Control (RTS/CTS) & Not required \\
  \hline
  no  & Modem Control Signals         & Not required \\
  \hline
\end{tabularx}
\end{table}

%------------------------
\subsubsection{HW Interfaces}
The following table lists all hardware interfaces relevant for normal operation.

\begin{table}[H]
\caption{HW Interfaces of Module UART0}
\label{tab:asUartPer02}
\centering
\begin{tabularx}{\textwidth}{|l |l |X|}
  \hline
   & Module internal name & Chip level name \\
  \hline
  \hline
  Supply &  &  \\
  \hline
  Clock & clk\_i & clk\_i \\
  \hline
  Bus & Wishbone slave & Wishbone B4 Classic \\
  \hline
  Interrupt & uart\_irq\_o & uart0\_irq\_s \\
  \hline
  DMA & None & - \\
  \hline
  External Signals & tx\_o & uart0\_tx\_o \\
                   & rx\_i & uart0\_rx\_i \\
  \hline
\end{tabularx}
\end{table}

%------------------------
\subsubsection{Bus Interface}
The UART controller exposes one bus interface:

\begin{description}
  \item[Wishbone Slave (configuration port)] Provides software access to all
    configuration registers, the TX/RX FIFOs, and the interrupt registers.
    Implemented by \texttt{as\_slave\_bpi}.
    Protocol: Wishbone B4 Classic, zero-wait-state, 64-bit data width.
    See Section~\ref{subsec:asSlaveBpi} (Architecture Concepts: Bus Concept).
\end{description}

%------------------------
\subsubsection{Registers}

\begin{table}[H]
\caption{Register List of Module UART0}
\label{tab:asUartPer03}
\centering
\begin{tabularx}{\textwidth}{|l |l |l |l |X|}
  \hline
  Offset & External Name & Internal Name & Mode & Description \\
  \hline
  \hline
  0x00 & \textbf{UART0\_ID}       & id\_reg\_s       & ro   & Peripheral identification \\
  \hline
  0x08 & \textbf{UART0\_LCR}      & lcr\_reg\_s      & rw   & Line Control: data bits, parity, stop bits, break \\
  \hline
  0x10 & \textbf{UART0\_CLKDIV}   & clkdiv\_reg\_s   & rw   & Baud rate divisor (16-bit) \\
  \hline
  0x18 & \textbf{UART0\_CTRL}     & ctrl\_reg\_s     & rw   & Control: loopback, tx\_flush, rx\_flush \\
  \hline
  0x20 & \textbf{UART0\_STATUS}   & status\_reg\_s   & ro   & TX/RX busy flags; sticky error flags \\
  \hline
  0x28 & \textbf{UART0\_DATA}     & data\_reg\_s     & rw   & Write: push TX FIFO; Read: pop RX FIFO \\
  \hline
  0x30 & \textbf{UART0\_FIFOSTAT} & fifostat\_reg\_s & ro   & TX/RX FIFO level, full, half, empty \\
  \hline
  0x38 & \textbf{UART0\_RXTHRES}  & rxthres\_reg\_s  & rw   & RX FIFO interrupt threshold (1\ldots16) \\
  \hline
  0x40 & \textbf{UART0\_RIS}      & ris\_reg\_s      & ro   & Raw Interrupt Status \\
  \hline
  0x48 & \textbf{UART0\_IMSC}     & imsc\_reg\_s     & rw   & Interrupt Mask Control \\
  \hline
  0x50 & \textbf{UART0\_MIS}      & mis\_reg\_s      & ro   & Masked Interrupt Status (= RIS\,\&\,IMSC) \\
  \hline
  0x58 & \textbf{UART0\_ICR}      & icr\_reg\_s      & wo   & Interrupt Clear (write-1-to-clear RIS bits) \\
  \hline
\end{tabularx}
\end{table}

%------------------------
\subsubsection{Related Sections of this Spec}
\begin{itemize}
  \item UART functional description starting with Section~\ref{subsec:uartfunc01}.
  \item Interrupt and DMA Concept in Chapter ``Architecture Concepts''.
  \item Peripheral Architecture Concept in Chapter ``Architecture Concepts'':
  \begin{itemize}
    \item Slave WB (Wishbone) BPI.
    \item Peripheral ID Registers.
    \item Peripheral Clocking Concept.
    \item Horizontal Interrupt and DMA Register Structure.
    \item Basic Interrupt and DMA Register Set.
    \item Synchronization Block.
    \item Standardized Flip-Flop based FIFO.
    \item TX-FIFO Registers.
    \item RX-FIFO Registers.
  \end{itemize}
  \item Requirements: Section~\ref{subsec:uartreq01}.
\end{itemize}

%=============================================================
% History of UART
%=============================================================
\subsection{History of UART}

\begin{table}[H]
\caption{History of UART Section}
\label{tab:uarthistory02}
\centering
\begin{tabularx}{\textwidth}{|X |X |X|}
  \hline
  \multicolumn{3}{|l|}{History} \\
  \hline
  \multicolumn{1}{|l}{Target Spec.} &
  \multicolumn{1}{l}{Current version: 0.1, 2026-06-05} &
  \multicolumn{1}{l|}{Author: A. Siggelkow} \\
  \multicolumn{1}{|l}{ } &
  \multicolumn{1}{l}{Previous version: none} &
  \multicolumn{1}{l|}{ } \\
  \hline
  \hline
  Paragraph (prev.) & Paragraph (current) / date & Changes \\
  \hline
  \hline
  \multicolumn{3}{|l|}{Chapter Version: C0.1; Date: 2026-06-05} \\
  \hline
  --- & \ref{sec:uartfunc01} / 2026-06-05 & First draft: Requirements + System Integration stubs (AS, Claude) \\
  \hline
\end{tabularx}
\end{table}

%=============================================================
% Functional Overview
%=============================================================
\subsection{Functional Overview}
\label{subsec:uartfunc01}
\textcolor{red}{To be written.}

The UART peripheral provides standard asynchronous serial communication in
half- or full-duplex mode.
It interfaces to the system via a Wishbone B4 Classic slave bus and
forwards received bytes to software via an RX FIFO and interrupts.

\paragraph{General Features:}
\begin{itemize}
  \item Configurable frame format (data bits, parity, stop bits)
  \item Programmable baud rate
  \item TX and RX FIFOs
  \item RX input synchroniser (metastability protection)
  \item Frame, parity, overrun, and break error detection
  \item RX timeout interrupt
  \item Loopback mode for self-test
  \item Interrupt generation (7 sources, maskable)
\end{itemize}

%=============================================================
% Structural Overview
%=============================================================
\subsection{Structural Overview}
\textcolor{red}{To be written.}

%\subsubsection{I/O of the UART Kernel}
%\subsubsection{I/O of the UART BPI}
%\subsubsection{Internals of the UART BPI}
%\subsubsection{I/O of the UART Top}

%=============================================================
% Service Requests
%=============================================================
\subsection{Service Requests of the UART}

Table~\ref{tab:asUartPer06b} shows the interrupt service requests of UART0.

\begin{table}[H]
\caption{UART0 Service Requests}
\label{tab:asUartPer06b}
\centering
\begin{tabularx}{\textwidth}{|l |l |l |X|}
  \hline
  Service Request & RIS Bit & IRQ & Explanation \\
  \hline
  \hline
  rx\_ready    & RIS[0] & uart0\_irq\_s & RX FIFO level $\geq$ RXTHRES. \\
  \hline
  tx\_ready    & RIS[1] & uart0\_irq\_s & TX FIFO level $<$ RXTHRES (space available). \\
  \hline
  rx\_timeout  & RIS[2] & uart0\_irq\_s & Data in RX FIFO; no new byte for $\geq$ 4 character times. \\
  \hline
  frame\_err   & RIS[3] & uart0\_irq\_s & Stop bit received as logic\;0. \\
  \hline
  parity\_err  & RIS[4] & uart0\_irq\_s & Received parity bit differs from computed parity. \\
  \hline
  overrun\_err & RIS[5] & uart0\_irq\_s & RX FIFO full; incoming byte discarded. \\
  \hline
  break\_det   & RIS[6] & uart0\_irq\_s & rx\_i held low for $\geq$ 1 full character time. \\
  \hline
\end{tabularx}
\end{table}

%=============================================================
% Peripheral Kernel
%=============================================================
\subsection{The UART Peripheral Kernel}
\label{subsec:uartfunc02}
\textcolor{red}{To be written.}

%=============================================================
% Registers
%=============================================================
\subsection{Registers}
\textcolor{red}{To be written -- detailed bit-field descriptions.}

\subsubsection{UART Register Overview}

\begin{table}[H]
\caption{UART Register List}
\label{tab:asUartRegList}
\centering
\begin{tabularx}{\textwidth}{|l |X |l |l|}
  \hline
  Register Group & Register Name & Symbol & Offset \\
  \hline
  \hline
  System Registers  & Identification Register     & \textbf{UART0\_ID}       & 0x00 \\
  \hline
  Configuration     & Line Control Register       & \textbf{UART0\_LCR}      & 0x08 \\
                    & Baud Rate Divisor Register  & \textbf{UART0\_CLKDIV}   & 0x10 \\
                    & Control Register            & \textbf{UART0\_CTRL}     & 0x18 \\
  \hline
  Status            & Status Register             & \textbf{UART0\_STATUS}   & 0x20 \\
                    & FIFO Status Register        & \textbf{UART0\_FIFOSTAT} & 0x30 \\
  \hline
  Data              & Data Register (TX/RX)       & \textbf{UART0\_DATA}     & 0x28 \\
  \hline
  Threshold         & RX Threshold Register       & \textbf{UART0\_RXTHRES}  & 0x38 \\
  \hline
  Interrupt         & Raw Interrupt Status        & \textbf{UART0\_RIS}      & 0x40 \\
  Registers         & Interrupt Mask Control      & \textbf{UART0\_IMSC}     & 0x48 \\
                    & Masked Interrupt Status     & \textbf{UART0\_MIS}      & 0x50 \\
                    & Interrupt Clear Register    & \textbf{UART0\_ICR}      & 0x58 \\
  \hline
\end{tabularx}
\end{table}

\subsubsection{Read-Write Features}
\begin{itemize}
  \item \textbf{w}: writable by SW
  \item \textbf{r}: readable by SW
  \item \textbf{h}: set by HW only
  \item \textbf{w1c}: write-1-to-clear
\end{itemize}

\subsubsection{LCR -- Line Control Register (offset 0x08, rw)}

\begin{table}[H]
\caption{UART0\_LCR Bit Fields}
\label{tab:asUartLCR}
\centering
\begin{tabularx}{\textwidth}{|l |l |l |X|}
  \hline
  Bits & Name & Mode & Description \\
  \hline
  \hline
  [1:0] & DATABITS & rw & Data word length:
    00\,=\,5\;bit, 01\,=\,6\;bit, 10\,=\,7\;bit, 11\,=\,8\;bit (reset). \\
  \hline
  [3:2] & PARITY   & rw & Parity mode:
    00\,=\,none (reset), 01\,=\,odd, 10\,=\,even, 11\,=\,reserved. \\
  \hline
  [4]   & STOPBITS & rw & Stop bit count:
    0\,=\,1 stop bit (reset), 1\,=\,2 stop bits. \\
  \hline
  [5]   & BREAK    & rw & Force break: 0\,=\,normal (reset), 1\,=\,TX held low. \\
  \hline
  [63:6] & --      & r  & Reserved; reads as 0. \\
  \hline
\end{tabularx}
\end{table}

\subsubsection{CTRL -- Control Register (offset 0x18, rw)}

\begin{table}[H]
\caption{UART0\_CTRL Bit Fields}
\label{tab:asUartCTRL}
\centering
\begin{tabularx}{\textwidth}{|l |l |l |X|}
  \hline
  Bits & Name & Mode & Description \\
  \hline
  \hline
  [0]    & LOOPBACK & rw & 0\,=\,normal (reset), 1\,=\,loopback (TX connected internally to RX). \\
  \hline
  [1]    & TX\_FLUSH & w  & Write\;1 to flush TX FIFO; auto-clears after one clock cycle. \\
  \hline
  [2]    & RX\_FLUSH & w  & Write\;1 to flush RX FIFO; auto-clears after one clock cycle. \\
  \hline
  [63:3] & --        & r  & Reserved; reads as 0. \\
  \hline
\end{tabularx}
\end{table}

\subsubsection{STATUS -- Status Register (offset 0x20, ro)}

\begin{table}[H]
\caption{UART0\_STATUS Bit Fields}
\label{tab:asUartSTATUS}
\centering
\begin{tabularx}{\textwidth}{|l |l |l |X|}
  \hline
  Bits & Name & Mode & Description \\
  \hline
  \hline
  [0]    & TX\_BUSY    & rh  & 1 while TX shift register is active. \\
  \hline
  [1]    & RX\_BUSY    & rh  & 1 while RX shift register is active (frame in progress). \\
  \hline
  [3:2]  & --          & r   & Reserved. \\
  \hline
  [4]    & FRAME\_ERR  & rh  & Sticky: frame error detected. Cleared by ICR[3]. \\
  \hline
  [5]    & PARITY\_ERR & rh  & Sticky: parity error detected. Cleared by ICR[4]. \\
  \hline
  [6]    & OVERRUN\_ERR & rh & Sticky: overrun error. Cleared by ICR[5]. \\
  \hline
  [7]    & BREAK\_DET  & rh  & Sticky: break condition detected. Cleared by ICR[6]. \\
  \hline
  [63:8] & --          & r   & Reserved; reads as 0. \\
  \hline
\end{tabularx}
\end{table}

\subsubsection{FIFOSTAT -- FIFO Status Register (offset 0x30, ro)}

\begin{table}[H]
\caption{UART0\_FIFOSTAT Bit Fields}
\label{tab:asUartFIFOSTAT}
\centering
\begin{tabularx}{\textwidth}{|l |l |l |X|}
  \hline
  Bits & Name & Mode & Description \\
  \hline
  \hline
  [4:0]  & TX\_LEVEL & rh & TX FIFO fill level (0\ldots FIFO\_DEPTH). \\
  \hline
  [5]    & TX\_EMPTY & rh & 1 when TX FIFO contains no entries. \\
  \hline
  [6]    & TX\_HALF  & rh & 1 when TX FIFO level\,$\leq$\,FIFO\_DEPTH/2. \\
  \hline
  [7]    & TX\_FULL  & rh & 1 when TX FIFO is full. \\
  \hline
  [15:8] & --        & r  & Reserved. \\
  \hline
  [20:16] & RX\_LEVEL & rh & RX FIFO fill level (0\ldots FIFO\_DEPTH). \\
  \hline
  [21]   & RX\_EMPTY & rh & 1 when RX FIFO contains no entries. \\
  \hline
  [22]   & RX\_HALF  & rh & 1 when RX FIFO level\,$\geq$\,FIFO\_DEPTH/2. \\
  \hline
  [23]   & RX\_FULL  & rh & 1 when RX FIFO is full. \\
  \hline
  [63:24] & --        & r  & Reserved; reads as 0. \\
  \hline
\end{tabularx}
\end{table}

\subsubsection{RIS -- Raw Interrupt Status (offset 0x40, ro)}

\begin{table}[H]
\caption{UART0\_RIS Bit Fields}
\label{tab:asUartRIS}
\centering
\begin{tabularx}{\textwidth}{|l |l |l |X|}
  \hline
  Bits & Name & Mode & Description \\
  \hline
  \hline
  [0]    & RX\_READY   & rh  & RX FIFO level\,$\geq$\,RXTHRES. \\
  \hline
  [1]    & TX\_READY   & rh  & TX FIFO level\,$<$\,RXTHRES (space available). \\
  \hline
  [2]    & RX\_TIMEOUT & rh  & Data in RX FIFO; no new byte for\,$\geq$\,4 character times. \\
  \hline
  [3]    & FRAME\_ERR  & rh  & Frame error (see UART-ER-001). Cleared by ICR[3]. \\
  \hline
  [4]    & PARITY\_ERR & rh  & Parity error (see UART-ER-002). Cleared by ICR[4]. \\
  \hline
  [5]    & OVERRUN\_ERR & rh & Overrun error (see UART-ER-003). Cleared by ICR[5]. \\
  \hline
  [6]    & BREAK\_DET  & rh  & Break detected (see UART-ER-004). Cleared by ICR[6]. \\
  \hline
  [63:7] & --           & r   & Reserved; reads as 0. \\
  \hline
\end{tabularx}
\end{table}

\subsubsection{ICR -- Interrupt Clear Register (offset 0x58, wo)}

\begin{table}[H]
\caption{UART0\_ICR Bit Fields}
\label{tab:asUartICR}
\centering
\begin{tabularx}{\textwidth}{|l |l |l |X|}
  \hline
  Bits & Name & Mode & Description \\
  \hline
  \hline
  [0]    & CLR\_RX\_READY    & w1c & Write\;1 clears RIS[0] (rx\_ready). \\
  \hline
  [1]    & CLR\_TX\_READY    & w1c & Write\;1 clears RIS[1] (tx\_ready). \\
  \hline
  [2]    & CLR\_RX\_TIMEOUT  & w1c & Write\;1 clears RIS[2] (rx\_timeout). \\
  \hline
  [3]    & CLR\_FRAME\_ERR   & w1c & Write\;1 clears RIS[3] and STATUS[4]. \\
  \hline
  [4]    & CLR\_PARITY\_ERR  & w1c & Write\;1 clears RIS[4] and STATUS[5]. \\
  \hline
  [5]    & CLR\_OVERRUN\_ERR & w1c & Write\;1 clears RIS[5] and STATUS[6]. \\
  \hline
  [6]    & CLR\_BREAK\_DET   & w1c & Write\;1 clears RIS[6] and STATUS[7]. \\
  \hline
  [63:7] & --                & --  & Reserved; writes ignored. \\
  \hline
\end{tabularx}
\end{table}
